\begin{ledgroupsized}[r]{120mm}
\footnotesize
\pstart
\noindent\textbf{Überlieferung:}
\pend
\end{ledgroupsized}
\begin{ledgroupsized}[r]{114mm}
\footnotesize
\pstart \parindent -6mm
\makebox[6mm][l]{\textit{L}}%
Aufzeichnung:
LH~XXXV~10, 7~Bl.~22. 
Ein Zettel (ca.~9~x~8,5~cm);
Ränder unregelmäßig beschnitten.
Zwei Seiten;
auf Bl.~22~r\textsuperscript{o} zwei nicht gestrichene Vorstufen zu \lbrack\textit{Fig.~1}\rbrack, die nicht wiedergegeben werden.
\pend
\end{ledgroupsized}
\vspace{5mm}
\begin{ledgroup}
\footnotesize
\pstart
\noindent%
\textbf{Datierungsgründe:} 
Bei der Berechnung der Geschwindigkeiten und Kräfte nach einem rechtwinkligen Stoß dreier Körper 
folgt Leibniz dem besonders in zwei Konzepten aus dem Zeitraum 1686 -- Oktober 1687 
ausführlich erarbeiteten Ansatz, d.\,h., der Dekomposition der ursprünglichen Kraft in die Katheten eines rechtwinkligen Dreiecks (siehe die Konzepte "Motuum compositioni non est fidendum, praesertim in concursu obliquo": \cite{02092}\textit{LSB} VIII,~3 N.\;65\textsubscript{3} und 
und "Principia universalia ad concursus directos ac obliquos determinandos": \cite{02093}N.\;66).
Wegen der sicheren Handhabung dieser Methode dürfte N.\;\ref{60048} nach den Stücken von 1686 entstanden sein.
\pend
\pstart
Mehrere ab 1690 verfasste Konzepte behandeln dieselbe Frage anhand einer neuen Methode: der Dekomposition des Dreikörpersystems in einfachere Teilsysteme, deren Kräfte separat berechnet und anschließend (als \textit{vires respectivae}) wieder zusammengesetzt werden können
(siehe bspw. N.\;\ref{60298}, N.\;\ref{58951}, N.\;\ref{58300}, N.\;\ref{58293}).
Dass Leibniz in N.\;\ref{60048} die ab 1686 entwickelte Methode verfolgt, die ab 1690 häufig auftretende hingegen nicht erwähnt, lässt die Vermutung zu, dass das Konzept spätestens um 1690 entstanden sei, wobei eine spätere Datierung nicht vollständig ausgeschlossen werden kann.
\pend
\end{ledgroup}
\selectlanguage{latin}
\vspace{8mm}
\pstart%
\normalsize%
\noindent%
\href{\Scan{LH_35_10_07_022r}}%
{\lbrack 22~r\textsuperscript{o}\rbrack}
\pend
\nobreak\begin{edfigure}[0.5em]%
\includegraphics[scale=1]{%
gesamttex/edit_VIII-4/images/LH_35_10_07_022_d_022r.pdf}
\edcaption{\edtext{\lbrack\textit{Fig.~1}\rbrack}{\lemma{}\Afootnote[nonum]{\textit{Neben \lbrack\textit{Fig.~1}\rbrack}: \textit{\ind{1}B} et \textit{\ind{2}B} hic coincidunt.}}}
\end{edfigure}
\pstart
\noindent
\edtext{$A=\lbrack D\rbrack=\lbrack E\rbrack$,}{%
\lemma{}%
\Bfootnote{%
$A = B = C$ \textit{L ändert Hrsg.}%
}}
\edtext{angulus \textit{\ind{3}D\ind{2}D}%
\textit{\ind{2}}\lbrack\textit{A}\rbrack\textit{\ind{2}E\ind{3}E}%
}{\lemma{angulus}%
\Bfootnote{%
\textit{(1)}~\textit{\ind{2}D\ind{2}C}\textit{\ind{2}E} %
\textit{(2)}~\textit{\ind{3}D\ind{2}D} %
\bars{\textit{\ind{2}C} \textit{ändert Hrsg.}} %
\textit{\ind{2}E\ind{3}E}~\textit{L}}}
rectus.%
\indexS{angulus rectus} \textit{C} centrum gravitatis commune,%
\indexS{centrum gravitatis commune} ejus via%
\indexS{via centri gravitatis} ante ictum\indexS{ictus} \textit{\ind{1}C\ind{2}C}, post ictum\indexS{ictus} 
\edtext{\textit{\ind{2}C\ind{3}}\lbrack\textit{D}\rbrack.}{%
\lemma{}%
\Bfootnote{%
\textit{\ind{2}C\ind{3}D} %
\textit{L ändert Hrsg.}}}
Debet esse $\ind{1}C\ind{2}C=\ind{2}C\ind{3}C$. $AC=\text{bis}\, BC$.  
\pend
\pstart
Quadr.\ $\ind{2}E\ind{3}E\,+\,\text{quadr.\ }\ind{2}D\ind{3}D$ est aequ.\ quater.\ quadr.\ \textit{\ind{2}B\ind{3}B}. 
\edtext{Igitur quadruplum}{%
\lemma{Igitur}%
\Bfootnote{%
\textit{(1)}~quadr.\ \textit{\ind{1}D\ind{2}D} %
\textit{(2)}~quadruplum~\textit{L}}}
quadratum \textit{\ind{1}B\ind{3}B} est %
vis\indexS{vis}
\edtext{communicata globis%
\indexS{globus}}{\lemma{communicata}\Bfootnote{%
\textit{(1)}~corporibus %
\textit{(2)}~globis~\textit{L}}}
\textit{D} et \textit{E}, cui si addatur quadratum \textit{\ind{2}A\ind{3}A}, habetur
\edtext{tota}{\lemma{}\Bfootnote{tota \textit{erg.~L}}}
vis\indexS{vis} aequalis priori, seu debet esse:
\pend
\pstart
4 $\overline{\text{quadr.}\ \textit{\ind{2}B\ind{3}B}}$ 
\edtext{+ qu.\ $\overline{\ind{2}A\ind{3}A}}{\lemma{}%
\Bfootnote{+ \bars{qu.\ \textit{erg.}} %
$\overline{\ind{2}A\ind{3}A}$~\textit{L}}}
\,=\,\text{qu.\ }\ind{1}A\ind{2}A$. 
\pend
\pstart
$\ind{1}A\ind{2}A=a$, $\ind{2}A\ind{2}E=\ind{2}A\ind{2}D=h$, et $\ind{2}A\ind{1}B=\overline{h:\sqrt2}$. 
Ergo \textit{\ind{1}A}\ind{1}$B=a+\overline{h:\sqrt2}$\lbrack,\rbrack\
$\ind{2}A\ind{2}C=\dfrac{2}{3}\ind{2}A\ind{1}B=h\sqrt2 :3$\lbrack,\rbrack\
$\ind{1}C\ind{1}A=\dfrac{2}{3}\ind{1}A\ind{1}B=2a:3+h\sqrt2 :3$. 
Jam $\ind{1}C\ind{2}A=\ind{2}A\ind{1}A-\ind{1}A\ind{1}C=a:3-h\sqrt2 :3$. 
Jam $\ind{1}C\ind{2}C=\ind{1}C\ind{2}A+\ind{2}A\ind{2}C$. 
Ergo $\ind{1}C\ind{2}C=a:3%
\edtext{=\ind{2}C\ind{3}C$.}{%
\lemma{}%
\Bfootnote{%
$=\ind{2}C\ind{3}C$. %
\textit{erg.~L}}}
Sit $\ind{1}B\ind{3}B=x$\lbrack,\rbrack\ adde
\edtext{$\ind{1}B\ind{2}C=h\sqrt2 :3$\lbrack,\rbrack}{%
\lemma{$\ind{1}B\ind{2}C=h\sqrt2 :3$}%
\Cfootnote{%
Der Wert von \textit{\ind{1}B\ind{2}C} ist eigentlich $h\sqrt2 :6$. Der Fehler wirkt sich auf die folgenden Ableitungen aus.}}
fit $\ind{3}B\ind{2}C=x+h\sqrt2 :3$,
unde aufer \textit{\ind{1}C\ind{2}C} seu $a:3$.
Fit $\ind{3}B\ind{3}C=x+h\sqrt2 :3-a:3$ cujus duplum est $\ind{3}A\ind{3}C=2x+2h\sqrt2 :3-2a:3$\lbrack,\rbrack\
aufer ab \textit{\ind{2}C\ind{3}C}\lbrack,\rbrack\ restat $a-2x-2h\sqrt2 :3=\ind{3}A\ind{2}C$\lbrack,\rbrack\ adde \textit{\ind{2}C\ind{2}A}\lbrack,\rbrack\ fit $a-2x+h\sqrt2 :3=\ind{2}A\ind{3}A$. \pend
\pstart
\href{\Scan{LH_35_10_07_022v}}%
{\lbrack 22~v\textsuperscript{o}\rbrack}\ 
Ergo fiet $4x^2 
\ovalbox{+a^2}
-4ax+2ah \sqrt2 :3+4x^2-4xh\sqrt2 :3+2h^2:9=\ovalbox{a^2}$
ob virium summam\indexS{summa virium}. 
\edtext{Sed si jam ponamus compendii causa \textit{h} esse infinite parvum, fit: $2x^2=ax$ seu 
\rule[-2.2ex]{0pt}{5.5ex}$x=\dfrac{1}{2}a$.}{%
\lemma{Sed \lbrack...\rbrack\  $x=\dfrac{1}{2}a$}%
\Cfootnote{%
Aufgrund der Setzung von \textit{h} als unendlich klein bleibt diese Ableitung von den Fehlern der vorangehenden Rechnungen unberührt und ist somit richtig.}}
Hoc est si puncta \textit{\ind{2}A\ind{2}D}\textit{\ind{2}E\ind{2}C} coincidere intelligantur seu globuli\indexS{globulus} sint infinite parvi fit \textit{\ind{2}B\ind{3}B} = dimidiae \textit{\ind{1}A\ind{2}A}, seu bis $\ind{2}B\ind{3}B=%
\edlabel{35_10_07_022_1a}%
\edtext{}{%
{\xxref%
{35_10_07_022_1a}{35_10_07_022_1b}}%
\lemma{\textit{\ind{1}A\ind{2}A}.}%
\Bfootnote{%
\textit{(1)}~Ergo quadrupla \textit{\ind{2}B\ind{3}B} = quad 
\textit{(2)}~Ergo~\textit{L}}}%
\ind{1}A\ind{2}A$. 
\pend
\pstart
Ergo\edlabel{35_10_07_022_1b} quadruplum quadratum \textit{\ind{2}B\ind{3}B} aequale
quadrato \textit{\ind{1}A\ind{2}A}. Jam supra ostensum est, quadruplum quadratum
\edtext{\textit{\ind{2}B\ind{3}B}, una}{\lemma{\textit{\ind{2}B\ind{3}B},}\Bfootnote{%
\textit{(1)}~unam %
\textit{(2)}~una~\textit{L}}}
cum quadrato \textit{\ind{2}A\ind{3}A}, aequali quadrato \textit{\ind{1}A\ind{2}A}. Ergo quadratum \textit{\ind{2}A\ind{3}A} aequatur nihilo, seu corpus \textit{A} post ictum\indexS{ictus} quiescit.  
\pend
\pstart
Si autem globuli\indexS{globulus} alicujus magnitudinis\indexS{magnitudo} ponantur, fit 
\pend
\pstart
\noindent
\rule[0mm]{0pt}{12mm}%
$x^2-\displaystyle\frac{1}{2}$%
{$ \left\{ \protect\begin{array}{l} a \\ h\sqrt2 :3 \end{array} \right. \!x$}
\rule[0mm]{0pt}{12mm}%
$\;+\;\displaystyle\frac{1}{4}$%
{$ \left\{ \begin{array}{l} a^2 \\ 2ah\sqrt2 :3 \\ 2h^2:3 \end{array} \right \}$}%
\rule[-4mm]{0pt}{12mm}%
$\;=\displaystyle\frac{1}{4}$%
{$\left \{\begin{array}{l} a^2 \\ 2ah\sqrt2 :3 \\2h^2:3 \end{array} \right \}$}%
$\begin{array}{l} \\ -2ah\sqrt2 :3 \\ -2h^2:9 \end{array}$
seu
\pend
\pstart
\noindent
 \rule[-4mm]{0pt}{12mm}%
 $x-\displaystyle\frac{1}{2}$%
 {$ \left \{\begin{array}{l} a \\ h\sqrt2 :3 \end{array} \right.$}%
 $ = \sqrt{a^2:4-ah:\sqrt2 -h^2:18}$. 
\pend
